\subsection{分次子模}

命题2. 设 A 是一个类型为 \(\Delta\) 的分次环，M 是一个类型为 \(\Delta\) 的分次 A-模， \((\mathbf{M}_{\lambda})\) 是它的分次，且 N 是 M 的一个子 A-模。以下性质等价：

\begin{enumerate}
\def\labelenumi{(\alph{enumi})}
\item
  \(\mathbf{N}\) 是族 \((\mathbf{N} \cap \mathbf{M}_{\lambda})_{\lambda \in \Delta}\) 之和。
\item
  N 的每个元素的齐次分量都属于 N。
\item
  N 由齐次元素生成。
\end{enumerate}

N 的每个元素都可以唯一地写成诸 \(M_{\lambda}\) 的元素之和，因而立即看出 (a) 与 (b) 等价且 (a) 蕴含 (c)。我们证明 (c) 蕴含 (b)。设 \((x_{t})_{t\in I}\) 是 N 的一族 \(\neq0\) 的齐次生成元，并令 \(\delta(t)\) 为 \(x_{t}\) 的次数。N 的每个元素都可写成 \(\sum_{i}a_{i}x_{i}\) ，其中 \(a_{i}\in A\) ；若 \(a_{i,\lambda}\) 是 \(a_{i}\) 的次数为 \(\lambda\) 的分量，则结论由关系式

\[
\sum_ {t \in I} \left(\sum_ {\mu \in \Delta} a _ {t, \mu} x _ {t}\right) = \sum_ {\lambda \in \Delta} \left(\sum_ {\mu + \delta (t) = \lambda} a _ {t, \mu} x _ {t}\right).
\]

注. (1) 在上述记号下，关系式 \(\sum_{t\in I}a_t x_t = 0\) 因而等价于关系式组 \(\sum_{\mu +\delta (t) = \lambda}a_{t,\mu}x_{t} = 0\) 。当 \(\Delta\) 是群时，这些关系式可写成 \(\sum_{t\in I}a_{t,\lambda -\delta (t)}x_{t} = 0\) 。

当 M 的子模 N 具有命题 2 中所述的等价性质时，显然 \(N \cap M_{\lambda}\) 构成与 N 的 A-模结构相容的一个分次，称为由 M 的分次诱导的分次；带有此分次的 N 称为 M 的分次子模。

推论 1. 若 N 是 M 的分次子模，且 \((x_{t})\) 是 N 的一个生成系，则 \(x_{t}\) 的齐次分量构成 N 的一个生成系。

推论 2. 若 N 是 M 的有限生成子模，则 N 具有由齐次元素组成的有限生成系。

只需应用推论 1，并注意到 M 的元素只有有限多个 \(\neq0\) 的齐次分量。

\(A_{s}\) （相应地 \(A_{a}\) ）的分次子模称为分次环 A 的分次左（相应地，右）理想。对 A 的每个子环 B，有 \(\mathrm{A}(\mathrm{B} \cap \mathrm{A}_{\lambda})(\mathrm{B} \cap \mathrm{A}_{\mu}) \subset \mathrm{B} \cap \mathrm{A}_{\lambda + \mu}\) ；若 B 是 A 的分次子 Z-模，则 A 的分次在 B 上诱导的分次便与 B 上的环结构相容；这时 B 称为 A 的分次子环。

显然，若 N（相应地，B）是 M 的分次子 A-模（相应地，A 的分次子环），则典范单射 \(N \rightarrow M\) （相应地 \(B \rightarrow A\) ）是 0 次分次模同态（相应地，分次环同态）。

若 N 是分次 A-模 M 的一个分次子模， \((\mathbf{M}_{\lambda})_{\lambda\in\Delta}\) 是 M 的分次，则 M/N 的子模 \((\mathbf{M}_{\lambda}+\mathbf{N})/\mathbf{N}\) 构成一个与该商模结构相容的分次。因为，若 \(N_{\lambda}=M_{\lambda}\cap N\) ，则 \((\mathbf{M}_{\lambda}+\mathbf{N})/\mathbf{N}\) 等同于 \(M_{\lambda}/N_{\lambda}\) ，且由命题 2 及 §1 第 6 号公式 (26) 可知，M/N 是它们的直和。此外，

\[
\mathrm{A} _ {\lambda} (\mathrm{M} _ {\mu} + \mathrm{N}) \subset \mathrm{A} _ {\lambda} \mathrm{M} _ {\mu} + \mathrm{N} \subset \mathrm{M} _ {\lambda + \mu} + \mathrm{N}
\]

，因此 \(\mathrm{A}_{\lambda}((\mathrm{M}_{\mu}+\mathrm{N})/\mathrm{N})\subset(\mathrm{M}_{\lambda+\mu}+\mathrm{N})/\mathrm{N}\) ，这就证明了我们的断言。分次 \(((\mathbf{M}_{\lambda} + \mathbf{N})/\mathbf{N})_{\lambda \in \Delta}\) 称为 M 的分次关于 N 的商分次，而带有此分次的商模 M/N 称为 M 关于分次子模 N 的分次商模；典范同态 \(M \to M/N\) 对此分次是次数为 0 的分次同态。

若 b 是 A 的分次双边理想，则 A/b 上的商分次与 A/b 上的环结构相容；带有此分次的环 A/b 称为 A 关于 b 的商分次环；典范同态 \(A \rightarrow A/b\) 对此分次是分次环的同态。

命题3. 设 A 是一个类型为 \(\Delta\) 的分次环，M、N 是两个类型为 \(\Delta\) 的分次 A-模，且 \(u\) : M \(\rightarrow\) N 是一个次数为 \(\delta\) 的分次 A-同态。那么：

\begin{enumerate}
\def\labelenumi{(\roman{enumi})}
\item
  \(\operatorname{Im}(u)\) 是 \(\mathbf{N}\) 的一个分次子模。
\item
  若 \(\delta\) 是 \(\Delta\) 的正则元素，则 \(\operatorname{Ker}(u)\) 是 \(\mathbf{M}\) 的一个分次子模。
\item
  若 \(\delta = 0\) ，则双射 \(\mathbf{M} / \mathrm{Ker}(u) \to \operatorname{Im}(u)\) （典范地对应于 \(u\) 的）是分次模的一个同构。
\end{enumerate}

断言 (i) 直接由定义和命题 2(c) 得出。若 \(x\) 是 \(\mathbf{M}\) 的一个元素，使得 \(u(x) = 0\) ，且 \(x = \sum_{\lambda} x_{\lambda}\) 是它的齐次分量分解（其中 \(x_{\lambda}\) 的次数为 \(\lambda\) ），则

\[
\sum_ {\lambda} u (x _ {\lambda}) = u (x) = 0
\]

且 \(u(x_{\lambda})\) 的次数为 \(\lambda + \delta\) ；若 \(\delta\) 正则，则关系 \(\lambda + \delta = \mu + \delta\) 蕴含 \(\lambda = \mu\) ，于是诸 \(u(x_{\lambda})\) 是 \(u(x)\) 的齐次分量，从而必然有 \(u(x_{\lambda}) = 0\) （对所有 \(\lambda \in \Delta\) ），这就证明了 (ii)。双射 \(v: \mathbf{M} / \operatorname{Ker}(u) \to \operatorname{Im}(u)\) （典范地对应于 \(u\) 的）这时是一个次数为 \(\delta\) 的分次同态，这由商分次的定义得出；由此得当 \(\delta = 0\) 时的 (iii)。

推论. 设 A、B 是两个类型为 \(\Delta\) 的分次环，且 u: A \(\rightarrow\) B 是分次环的一个分次同态。则 \(\operatorname{Im}(u)\) 是 B 的一个分次子模， \(\operatorname{Ker}(u)\) 是 A 的一个分次双边理想，而典范地对应于 u 的双射 A/Ker(u) \(\rightarrow\) \(\operatorname{Im}(u)\) 是分次环的一个同构。

只需将命题 3 应用于将 u 视为分次 Z-模的 0 次同态即可。

命题4. 设 A 是一个类型为 \(\Delta\) 的分次环，M 是一个类型为 \(\Delta\) 的分次 A-模。

\begin{enumerate}
\def\labelenumi{(\roman{enumi})}
\item
  \(\mathbf{M}\) 的分次子模的每一和与每一交都是一个分次子模。
\item
  若 \(x\) 是 \(\mathbf{M}\) 的一个次数为 \(\mu\) （在 \(\Delta\) 中可消去）的齐次元素，则 \(x\) 的零化子是 \(\mathbf{A}\) 的一个分次左理想。
\item
  若 \(\Delta\) 的所有元素都可消去，则 M 的一个分次子模的零化子是 A 的一个分次双边理想。
\end{enumerate}

若 \((\mathbf{N}_{t})\) 是 M 的一族分次子模，则命题 2 的性质 (c) 表明诸 \(N_{t}\) 之和由齐次元素生成，而命题 2 的性质 (b) 证明 \(\bigcap_{t} N_{t}\) 中每个元素的齐次分量都属于 \(\bigcap_{t} N_{t}\) ；由此得 (i)。

为证明 (ii)，只需注意到 \(\operatorname{Ann}(x)\) 是同态 \(a \mapsto ax\) （ A-模 \(\mathbf{A}_s\) 到 \(\mathbf{M}\) 中的）的核，且该同态是次数为 \(\mu\) 的分次同态；结论由命题 3(ii) 得出。最后，(iii) 是 (i) 与 (ii) 的推论，因为分次子模 \(\mathbf{N}\) （ \(\mathbf{M}\) 的）的零化子是 \(\mathbf{N}\) 的齐次元素的零化子之交（据命题 2）。

注2. 设 M 为分次 A-模，E 为 M 的子模；由命题4(i)可知，存在包含于 E 中的 M 的最大分次子模 N'，以及包含 E 的 M 的最小分次子模 N''；N' 是 E 中所有其齐次分量均属于 E 的元素 x 的集合，而 N'' 是由 E 的一个生成系的齐次分量所生成的 M 的子模。

命题5. 设 A 是一个类型为 \(\Delta\) 的分次环。若 \(\Delta\) 的每个元素都可消去，则对每个齐次元素 \(a \in \mathbf{A}\) ， \(a\) 在 A 中的中心化子（I, §1, 第 5 号）是 A 的一个分次子环。

设 \(a\) 的次数为 \(\delta\) ；令 \(b = \sum_{\lambda} b_{\lambda}\) 是一个与 \(a, b_{\lambda}\) 可交换的元素，后者是 \(b\) 的次数为 \(\lambda\) 的齐次分量（对所有 \(\lambda \in \Delta\) ）。于是按假设 \(\sum_{\lambda} (ab_{\lambda} - b_{\lambda}a) = 0\) ，且 \(ab_{\lambda} - b_{\lambda}a\) 是次数为 \(\lambda + \delta\) 的齐次元素；由于 \(\delta\) 可消去，由此得出 \(ab_{\lambda} = b_{\lambda}a\) （对所有 \(\lambda\) ），这就证明了我们的断言。

推论. 若 \(\Delta\) 的每个元素都可消去，则 A 的分次子环 B 的中心化子（特别地，A 的中心）是 A 的一个分次子环。

它是 B 的齐次元素的中心化子的交集。

注. (3) 分次环的一个直系统 \((\mathbf{A}_{\alpha}, \phi_{\beta\alpha})\) （类型为 \(\Delta\) 的）（相应地，直系统 \((\mathbf{M}_{\alpha}, f_{\beta\alpha})\) ——分次 \(A_{\alpha}\) -模的，类型为 \(\Delta\) 的）是指环（相应地， \(A_{\alpha}\) -模）的一个直系统，使得每个 \(A_{\alpha}\) （相应地， \(M_{\alpha}\) ）是类型为 \(\Delta\) 的分次的，且每个 \(\phi_{\beta\alpha}\) （相应地， \(f_{\beta\alpha}\) ）是分次环的同态（相应地，分次模的次数为 0 的 \(A_{\alpha}\) -同态）。若 \((\mathbf{A}_{\alpha}^{\lambda})_{\lambda \in \Delta}\) （相应地， \((\mathbf{M}_{\alpha}^{\lambda})_{\lambda \in \Delta}\) ）是 \(A_{\alpha}\) （相应地， \(M_{\alpha}\) ）的分次，且我们记

\[
\mathrm{A} = \lim _ {\rightarrow} \mathrm{A} _ {\alpha}, \quad \mathrm{A} ^ {\lambda} = \lim _ {\overrightarrow {\alpha}} \mathrm{A} _ {\alpha} ^ {\lambda} \quad (\text { resp.   } \mathrm{M} = \lim _ {\rightarrow} \mathrm{M} _ {\alpha}, \mathrm{M} ^ {\lambda} = \lim _ {\overrightarrow {\alpha}} \mathrm{M} _ {\alpha} ^ {\lambda}),
\]

由 § 6 第 2 号命题 5 可知， \((\mathbf{A}^{\wedge})\) （相应地， \((\mathbf{M}^{\wedge}))\) 是 A（相应地，M）的一个分次，且由 I, §10, 第 3 号与第 4 号可知，该分次与 A 上的环结构（相应地，M 上的 A-模结构）相容。分次环 A（相应地，分次 A-模 M）称为分次环的直系统 \((\mathbf{A}_{\alpha}, \phi_{\beta \alpha})\) （相应地，分次模的直系统 \((\mathbf{M}_{\alpha}, f_{\beta \alpha})\) ）的直极限。若 \(\phi_{\alpha} \colon \mathbf{A}_{\alpha} \to \mathbf{A}\) （相应地， \(f_{\alpha} \colon \mathbf{M}_{\alpha} \to \mathbf{M}\) ）是典范映射，则 \(\phi_{\alpha}\) （相应地， \(f_{\alpha}\) ）是分次环的同态（相应地，分次 \(\mathbf{A}_{\alpha}\) -模的次数为 0 的同态）。

